% !TeX program = lualatex
% Template no oficial para presentaciones de posgrado en Ingeniería UC.
% Edita los comandos de esta sección para adaptar el template a Magíster o Doctorado.
\documentclass{ucbeamer}

\addbibresource{references.bib}

% ---------------------------------------------------------
% Datos globales editables de la presentación
% ---------------------------------------------------------
\newcommand{\tipoPrograma}{Magíster / Doctorado}
\newcommand{\nombrePrograma}{Nombre del programa de posgrado}
\newcommand{\tituloPresentacion}{Título de la presentación}
\newcommand{\tituloTesis}{Título tentativo de la tesis}
\newcommand{\nombreEstudiante}{Nombre del estudiante}
\newcommand{\nombreGuia}{Nombre profesor/a guía}
\newcommand{\nombreCoGuia}{Nombre profesor/a co-guía}
\newcommand{\departamentoEscuela}{Departamento / Escuela}
\newcommand{\fechaPresentacion}{Santiago, Chile -- Fecha}
\newcommand{\correoInstitucional}{correo.institucional@uc.cl}

% ---------------------------------------------------------
% Portada y metadatos Beamer
% ---------------------------------------------------------
\title{\bf \tituloPresentacion}
\subtitle{Presentación de avance de posgrado}
\author[\nombreEstudiante]{\nombreEstudiante}
\date[Fecha]{\fechaPresentacion}
\institute{Pontificia Universidad Católica de Chile}

\setprogramtype{\tipoPrograma}
\setthesistype{Tesis de \tipoPrograma: \tituloTesis}
\setprogram{\nombrePrograma}
\setdepartment{\departamentoEscuela}
\setfaculty{Escuela de Ingeniería}
\setsupervisor{Profesor/a guía: \nombreGuia}
\setcosupervisor{Profesor/a co-guía: \nombreCoGuia}
\setcontact{\correoInstitucional}
\setfootertext{Presentación de posgrado -- Ingeniería UC}

% Para reemplazar logos:
% \UCsetlogos{assets/logos/uc_logo_template.png}{assets/logos/department_logo_placeholder.png}

\begin{document}

\UCtitleframe

\begin{frame}{Ruta de la presentación}
  \begin{columns}[T,onlytextwidth]
    \begin{column}{0.56\textwidth}
      \tableofcontents
    \end{column}
    \begin{column}{0.40\textwidth}
      \begin{UCKeyPoint}
        Template no oficial para organizar una presentación de avance de posgrado. Reemplaza cada bloque por contenido propio y elimina las láminas que no correspondan a tu programa.
      \end{UCKeyPoint}
      \UCmetric{16:9}{Formato}{Portada, contexto, problema, objetivos, metodología, avance, resultados, cronograma, referencias y anexos.}
    \end{column}
  \end{columns}
\end{frame}

% =========================================================
\section{Contexto y problema}
% =========================================================

\begin{frame}{Contexto general}
  \begin{columns}[T,onlytextwidth]
    \begin{column}{0.50\textwidth}
      \begin{UCBlock}{Idea central}
        Inserte aquí el contexto general del área de investigación, la relevancia académica o aplicada del tema y el escenario donde se sitúa la tesis de posgrado.
      \end{UCBlock}
      \begin{itemize}
        \item Antecedente principal del problema.
        \item Brecha o necesidad que motiva el estudio.
        \item Alcance general del proyecto de Magíster/Doctorado.
      \end{itemize}
    \end{column}
    \begin{column}{0.46\textwidth}
      \centering
      \begin{tikzpicture}[scale=0.92,every node/.style={font=\small}]
        \fill[UCBlueLight] (0,0) rectangle (6.3,4.15);
        \draw[UCBluePPT,line width=0.8pt] (0,0) rectangle (6.3,4.15);
        \node[draw=UCBluePPT,fill=UCWhite,rounded corners=2pt,align=center,text width=1.85cm] (a) at (1.15,3.25) {Variable\\de entrada};
        \node[draw=UCBluePPT,fill=UCWhite,rounded corners=2pt,align=center,text width=1.72cm] (b) at (3.25,3.25) {Proceso\\o método};
        \node[draw=UCBluePPT,fill=UCWhite,rounded corners=2pt,align=center,text width=1.70cm] (c) at (5.25,3.25) {Condición\\de análisis};
        \node[draw=UCYellowDeep,fill=UCYellow!18!white,rounded corners=2pt,align=center,text width=4.25cm] (d) at (3.15,1.65) {Sistema de estudio\\o caso de aplicación};
        \node[draw=UCBlueDark,fill=UCBlueDark,text=UCWhite,rounded corners=2pt,align=center,text width=2.55cm] (e) at (3.15,0.42) {Resultado\\de ejemplo};
        \draw[-{Latex[length=2mm]},UCBluePPT,line width=0.8pt] (a) -- (d);
        \draw[-{Latex[length=2mm]},UCBluePPT,line width=0.8pt] (b) -- (d);
        \draw[-{Latex[length=2mm]},UCBluePPT,line width=0.8pt] (c) -- (d);
        \draw[-{Latex[length=2mm]},UCYellowDeep,line width=0.9pt] (d) -- (e);
      \end{tikzpicture}
      {\scriptsize Figura de ejemplo: reemplazar por esquema conceptual propio.}
    \end{column}
  \end{columns}
\end{frame}

\begin{frame}{Motivación}
  \begin{columns}[T,onlytextwidth]
    \begin{column}{0.48\textwidth}
      \begin{UCBlock}{Motivación académica}
        Inserte aquí por qué el tema es relevante para la literatura, la disciplina o la formación de conocimiento en ingeniería.
      \end{UCBlock}
      \begin{UCNeutralBox}
        Consejo: conecta la motivación con una brecha verificable, no solo con una descripción general del área.
      \end{UCNeutralBox}
    \end{column}
    \begin{column}{0.48\textwidth}
      \begin{UCBlock}{Motivación aplicada}
        Inserte aquí la relevancia práctica, industrial, social, tecnológica o metodológica del proyecto.
      \end{UCBlock}
      \UCsmallmetric{Impacto}{Variable de ejemplo}{Indicar el aporte esperado en una frase breve.}
      \UCsmallmetric{Escala}{Caso de ejemplo}{Indicar si el estudio es conceptual, experimental, computacional o aplicado.}
    \end{column}
  \end{columns}
\end{frame}

\begin{frame}{Problema de investigación}
  \begin{columns}[T,onlytextwidth]
    \begin{column}{0.52\textwidth}
      \begin{UCBlock}{Problema}
        Inserte aquí el problema de investigación. Debe quedar claro qué no se conoce, qué no está resuelto o qué requiere ser evaluado.
      \end{UCBlock}
      \begin{itemize}
        \item Elemento 1 del problema.
        \item Elemento 2 del problema.
        \item Elemento 3 del problema.
      \end{itemize}
    \end{column}
    \begin{column}{0.44\textwidth}
      \begin{UCKeyPoint}
        \textbf{Mensaje clave.} La lámina debe permitir que la comisión entienda por qué la tesis de posgrado es necesaria y qué decisión metodológica se deriva del problema.
      \end{UCKeyPoint}
    \end{column}
  \end{columns}
\end{frame}

\begin{frame}{Pregunta de investigación}
  \begin{columns}[T,onlytextwidth]
    \begin{column}{0.52\textwidth}
      \begin{UCBlock}{Pregunta principal}
        ¿Cómo influye \textit{variable de ejemplo} sobre \textit{resultado de ejemplo} en \textit{sistema o caso de estudio} bajo \textit{condiciones de análisis}?
      \end{UCBlock}
    \end{column}
    \begin{column}{0.44\textwidth}
      \begin{UCKeyPoint}
        \textbf{Hipótesis o supuesto de trabajo.} Inserte aquí la hipótesis, con lenguaje prudente y verificable.
      \end{UCKeyPoint}
    \end{column}
  \end{columns}
\end{frame}

% =========================================================
\section{Objetivos y antecedentes}
% =========================================================

\begin{frame}{Objetivo general y objetivos específicos}
  \begin{UCBlock}{Objetivo general}
    Inserte aquí el objetivo general de la tesis de posgrado, formulado con un verbo claro y alineado con la pregunta de investigación.
  \end{UCBlock}
  \vspace{0.08cm}
  \begin{enumerate}
    \item Inserte aquí el primer objetivo específico.
    \item Inserte aquí el segundo objetivo específico.
    \item Inserte aquí el tercer objetivo específico.
    \item Inserte aquí el cuarto objetivo específico, si corresponde.
  \end{enumerate}
\end{frame}

\begin{frame}{Marco teórico o antecedentes}
  \begin{columns}[T,onlytextwidth]
    \begin{column}{0.32\textwidth}
      \begin{UCBlock}{Concepto 1}
        Inserte aquí una definición, teoría o referencia clave.
      \end{UCBlock}
    \end{column}
    \begin{column}{0.32\textwidth}
      \begin{UCBlock}{Concepto 2}
        Inserte aquí un antecedente metodológico o técnico relevante.
      \end{UCBlock}
    \end{column}
    \begin{column}{0.32\textwidth}
      \begin{UCBlock}{Concepto 3}
        Inserte aquí el antecedente que conecta con la brecha.
      \end{UCBlock}
    \end{column}
  \end{columns}
  \begin{UCKeyPoint}
    Evita convertir esta sección en una revisión bibliográfica extensa. Prioriza los conceptos que justifican el método y los objetivos.
  \end{UCKeyPoint}
\end{frame}

\begin{frame}{Alcance de la tesis de posgrado}
  \begin{columns}[T,onlytextwidth]
    \begin{column}{0.58\textwidth}
      \centering
      \begin{tikzpicture}[x=1cm,y=1cm]
        \node (a) at (0,3.12) {\UCProcessArrow{Entrada}{datos, muestras o caso}};
        \node (b) at (2.05,3.12) {\UCProcessArrow{Preparación}{limpieza o diseño}};
        \node (c) at (4.42,3.12) {\UCProcessArrow{Análisis}{modelo o ensayo}};
        \node (d) at (0.70,1.36) {\UCProcessArrow{Validación}{criterio definido}};
        \node (e) at (2.94,1.36) {\UCProcessArrow{Resultados}{indicadores clave}};
        \node (f) at (5.16,1.36) {\UCProcessArrow{Cierre}{discusión y límites}};
        \draw[-{Latex[length=2mm]},UCBluePPT,line width=0.7pt] (0.95,3.12)--(1.56,3.12);
        \draw[-{Latex[length=2mm]},UCBluePPT,line width=0.7pt] (3.02,3.12)--(3.78,3.12);
        \draw[-{Latex[length=2mm]},UCBluePPT,line width=0.7pt] (4.40,2.62)--(0.75,1.89);
        \draw[-{Latex[length=2mm]},UCBluePPT,line width=0.7pt] (4.40,2.62)--(2.92,1.89);
        \draw[-{Latex[length=2mm]},UCYellowDeep,line width=0.8pt] (3.62,1.36)--(4.22,1.36);
      \end{tikzpicture}
    \end{column}
    \begin{column}{0.38\textwidth}
      \begin{UCKeyPoint}
        Indica explícitamente qué queda dentro y fuera del alcance. Esto ayuda a controlar expectativas durante el avance.
      \end{UCKeyPoint}
    \end{column}
  \end{columns}
\end{frame}

% =========================================================
\section{Metodología}
% =========================================================

\begin{frame}{Diseño metodológico}
  \begin{table}
  \fontsize{7.9}{9.2}\selectfont
  \begin{tabularx}{\textwidth}{>{\bfseries}l X X X}
    \toprule
    Etapa & Insumo principal & Método de ejemplo & Producto esperado \\
    \midrule
    1 & Datos o antecedentes & Revisión, levantamiento o caracterización & Base de análisis \\
    2 & Variables seleccionadas & Modelo, experimento o algoritmo & Resultado preliminar \\
    3 & Caso de estudio & Validación, sensibilidad o comparación & Evidencia principal \\
    4 & Resultados integrados & Discusión técnica & Conclusiones y próximos pasos \\
    \bottomrule
  \end{tabularx}
  \end{table}
  \begin{UCKeyPoint}
    Reemplaza esta tabla por tu metodología real. La estructura privilegia trazabilidad: insumo, método y producto.
  \end{UCKeyPoint}
\end{frame}

\begin{frame}{Variables, indicadores y criterios}
  \begin{columns}[T,onlytextwidth]
    \begin{column}{0.33\textwidth}
      \UCmetric{X}{Variable de entrada}{Describir unidad, rango o fuente.}
      \UCmetric{Y}{Resultado de ejemplo}{Describir indicador principal.}
    \end{column}
    \begin{column}{0.33\textwidth}
      \UCmetric{$n$}{Muestras / casos}{Indicar tamaño, cobertura o alcance.}
      \UCmetric{$\Delta$}{Criterio de cambio}{Indicar umbral o comparación.}
    \end{column}
    \begin{column}{0.30\textwidth}
      \begin{UCBlock}{Control de calidad}
        Inserte aquí criterios de reproducibilidad, validación, supuestos o limitaciones metodológicas.
      \end{UCBlock}
    \end{column}
  \end{columns}
\end{frame}

\begin{frame}{Flujo de trabajo}
  \begin{columns}[T,onlytextwidth]
    \begin{column}{0.24\textwidth}
      \centering
      \UCProcessArrow{1. Datos}{fuentes, muestras o información}
      \vspace{0.30cm}
      \begin{tikzpicture}[scale=0.54]
        \fill[UCBlueLight] (0,0) rectangle (2.6,1.2);
        \draw[UCBluePPT] (0,0)--(2.6,1.2);
        \draw[UCBluePPT] (0,1.2)--(2.6,0);
      \end{tikzpicture}
    \end{column}
    \begin{column}{0.24\textwidth}
      \centering
      \UCProcessArrow{2. Método}{modelo, experimento o análisis}
      \vspace{0.42cm}
      \begin{tikzpicture}[scale=0.54]
        \foreach \x in {0,0.75,1.5}{\foreach \y in {0,0.55,1.1}{\draw[UCBluePPT,fill=UCWhite] (\x,\y) circle (0.20);}}
        \fill[UCYellow] (1.5,1.1) circle (0.20);
      \end{tikzpicture}
    \end{column}
    \begin{column}{0.24\textwidth}
      \centering
      \UCProcessArrow{3. Evaluación}{indicadores y sensibilidad}
      \vspace{0.30cm}
      \begin{tikzpicture}[scale=0.54]
        \draw[UCBluePPT,thick] (0.2,0.2)--(0.7,1.1)--(1.2,0.65)--(1.8,1.5)--(2.2,0.3);
        \draw[UCGray] (0.15,0.15)--(2.3,0.15);
      \end{tikzpicture}
    \end{column}
    \begin{column}{0.24\textwidth}
      \centering
      \UCProcessArrow{4. Discusión}{resultados y límites}
      \vspace{0.30cm}
      \begin{tikzpicture}[scale=0.54]
        \draw[UCBluePPT,thick] (0,0) -- (2.4,0) -- (2.4,1.5) -- (0,1.5) -- cycle;
        \draw[UCYellowDeep,line width=1pt] (0.25,0.3)--(0.9,0.7)--(1.35,0.55)--(2.1,1.15);
      \end{tikzpicture}
    \end{column}
  \end{columns}
  \vspace{0.18cm}
  \begin{UCKeyPoint}
    Usa esta lámina para explicar el método completo antes de entrar al detalle técnico.
  \end{UCKeyPoint}
\end{frame}

% =========================================================
\section{Avance y resultados preliminares}
% =========================================================

\begin{frame}{Estado actual del avance}
  \begin{columns}[T,onlytextwidth]
    \begin{column}{0.32\textwidth}
      \UCsmallmetric{Completado}{Etapa 1}{Inserte avance realizado.}
      \UCsmallmetric{En curso}{Etapa 2}{Inserte trabajo actual.}
    \end{column}
    \begin{column}{0.32\textwidth}
      \UCsmallmetric{Pendiente}{Etapa 3}{Inserte tarea pendiente.}
      \UCsmallmetric{Riesgo}{Etapa crítica}{Inserte riesgo y mitigación.}
    \end{column}
    \begin{column}{0.32\textwidth}
      \begin{UCBlock}{Lectura general}
        Inserte aquí un resumen breve del estado del proyecto, destacando avances verificables y decisiones pendientes.
      \end{UCBlock}
    \end{column}
  \end{columns}
\end{frame}

\begin{frame}{Plantilla para resultado principal}
  \begin{columns}[T,onlytextwidth]
    \begin{column}{0.58\textwidth}
      \begin{tikzpicture}
        \begin{axis}[
          ybar, width=\linewidth, height=5.0cm,
          ymin=0, ymax=100,
          ylabel={Indicador de ejemplo, \%},
          symbolic x coords={A,B,C,D,E}, xtick=data,
          grid=major, grid style={UCGrayLight},
          axis line style={UCBlueDark},
          tick label style={font=\scriptsize}, label style={font=\scriptsize},
          bar width=12pt,
          nodes near coords, nodes near coords style={font=\scriptsize, color=UCBlueDark},
        ]
        \addplot+[fill=UCBluePPT,draw=UCBluePPT] coordinates {(A,18)(B,35)(C,42)(D,55)(E,61)};
        \end{axis}
      \end{tikzpicture}
    \end{column}
    \begin{column}{0.38\textwidth}
      \UCmetric{+XX}{Resultado de ejemplo}{Comparación contra línea base.}
      \UCmetric{$Y$}{Variable explicativa}{Conectar con hipótesis o método.}
      \begin{UCKeyPoint}
        Mensaje de una frase: interpretación principal del resultado preliminar.
      \end{UCKeyPoint}
    \end{column}
  \end{columns}
\end{frame}

\begin{frame}{Plantilla para comparar tendencias}
  \begin{columns}[T,onlytextwidth]
    \begin{column}{0.64\textwidth}
      \begin{tikzpicture}
        \begin{axis}[
          width=\linewidth,height=5.15cm,
          xlabel={Tiempo o condición}, ylabel={Variable de respuesta},
          xmin=0,xmax=10,ymin=0,ymax=1.0,
          grid=both, grid style={UCGrayLight},
          axis line style={UCBlueDark},
          tick label style={font=\scriptsize}, label style={font=\scriptsize},
          legend columns=1,
          legend style={font=\scriptsize,draw=none,fill=none,at={(0.03,0.97)},anchor=north west}
        ]
        \addplot+[UCGray,mark=triangle*,line width=1pt] table[x=x,y=baseline,col sep=comma] {assets/data/example_results.csv};
        \addplot+[UCBluePPT,mark=*,line width=1pt] table[x=x,y=scenario_a,col sep=comma] {assets/data/example_results.csv};
        \addplot+[UCYellowDeep,mark=square*,line width=1pt] table[x=x,y=scenario_b,col sep=comma] {assets/data/example_results.csv};
        \legend{Base, Escenario A, Escenario B}
        \end{axis}
      \end{tikzpicture}
    \end{column}
    \begin{column}{0.32\textwidth}
      \begin{UCBlock}{Cómo narrarla}
        \begin{itemize}
          \item Tendencia principal.
          \item Diferencia entre escenarios.
          \item Incertidumbre o limitación.
          \item Lectura para próximos pasos.
        \end{itemize}
      \end{UCBlock}
    \end{column}
  \end{columns}
\end{frame}

\begin{frame}{Discusión preliminar}
  \begin{columns}[T,onlytextwidth]
    \begin{column}{0.31\textwidth}
      \begin{UCBlock}{Dato}
        Resultado medido, simulado, observado o calculado.
      \end{UCBlock}
    \end{column}
    \begin{column}{0.31\textwidth}
      \begin{UCBlock}{Evidencia}
        Tendencia reproducible, contraste o patrón relevante.
      \end{UCBlock}
    \end{column}
    \begin{column}{0.31\textwidth}
      \begin{UCBlock}{Interpretación}
        Explicación técnica respaldada por los resultados.
      \end{UCBlock}
    \end{column}
  \end{columns}
  \begin{UCKeyPoint}
    Evita presentar conclusiones definitivas si el avance todavía es parcial. Usa lenguaje como “sugiere”, “es consistente con” o “requiere verificación”.
  \end{UCKeyPoint}
\end{frame}

% =========================================================
\section{Plan de trabajo y cierre}
% =========================================================

\begin{frame}{Cronograma}
  \begin{table}
  \fontsize{7.8}{9.0}\selectfont
  \begin{tabularx}{\textwidth}{>{\bfseries}l *{6}{>{\centering\arraybackslash}X}}
    \toprule
    Actividad & M1 & M2 & M3 & M4 & M5 & M6 \\
    \midrule
    Revisión y ajuste metodológico & \cellcolor{UCBlueSoft} & \cellcolor{UCBlueSoft} & & & & \\
    Desarrollo / experimentación & & \cellcolor{UCBlueSoft} & \cellcolor{UCBlueSoft} & \cellcolor{UCBlueSoft} & & \\
    Análisis de resultados & & & \cellcolor{UCYellow!35} & \cellcolor{UCYellow!35} & \cellcolor{UCYellow!35} & \\
    Escritura y defensa & & & & & \cellcolor{UCBlueSoft} & \cellcolor{UCBlueSoft} \\
    \bottomrule
  \end{tabularx}
  \end{table}
  \begin{UCKeyPoint}
    Reemplaza este cronograma por tu planificación real. Usa meses relativos o fechas concretas según lo que exija el seminario.
  \end{UCKeyPoint}
\end{frame}

\begin{frame}{Próximos pasos}
  \begin{columns}[T,onlytextwidth]
    \begin{column}{0.50\textwidth}
      \begin{UCBlock}{Trabajo técnico}
        \begin{itemize}
          \item Completar actividad pendiente 1.
          \item Completar actividad pendiente 2.
          \item Validar resultado o supuesto crítico.
        \end{itemize}
      \end{UCBlock}
    \end{column}
    \begin{column}{0.46\textwidth}
      \begin{UCBlock}{Trabajo de tesis de posgrado}
        \begin{itemize}
          \item Actualizar marco teórico.
          \item Consolidar resultados preliminares.
          \item Preparar capítulo o informe de avance.
        \end{itemize}
      \end{UCBlock}
    \end{column}
  \end{columns}
\end{frame}

\begin{frame}{Cierre}
  \begin{enumerate}
    \item \textbf{Contexto:} recordar el problema y su relevancia.
    \item \textbf{Avance:} resumir qué se completó y qué evidencia existe.
    \item \textbf{Limitaciones:} declarar supuestos, restricciones o incertidumbres.
    \item \textbf{Proyección:} indicar próximos pasos y decisiones esperadas.
  \end{enumerate}
  \begin{UCKeyPoint}
    Consejo: cerrar con tres o cuatro mensajes, cada uno conectado con una lámina anterior.
  \end{UCKeyPoint}
\end{frame}

\section{Referencias}

\begin{frame}[allowframebreaks]{Referencias}
  \nocite{*}
  \printbibliography
\end{frame}

\UCthanksframe[Gracias]

\appendix
\section{Anexos}

\begin{frame}{Anexo: respaldo técnico}
  \begin{itemize}
    \item Tabla de ejemplo con datos complementarios.
    \item Figura de ejemplo con validación o sensibilidad.
    \item Detalle metodológico que no cabe en la presentación principal.
    \item Referencias adicionales o supuestos utilizados.
  \end{itemize}
\end{frame}


\end{document}
